% Standalone TikZ figure: the piecewise-constant prediction surface of a decision tree.
% \input from the chapter's raw-LaTeX block; not rendered by Obsidian.
\begin{figure}[H]
\centering
\begin{tikzpicture}[scale=1.0, >=stealth]
    \definecolor{regiongray}{HTML}{BDBDBD}

    % Oblique projection: (x1,x2,prediction) -> (x1 + 0.45*x2, 0.32*x2 + 0.85*prediction).
    % Base plane (the feature space at prediction zero).
    \fill[black!3]
        ({0+0.45*0},{0.32*0}) -- ({4+0.45*0},{0.32*0}) --
        ({4+0.45*3},{0.32*3}) -- ({0+0.45*3},{0.32*3}) -- cycle;
    \draw[black!25]
        ({0+0.45*0},{0.32*0}) -- ({4+0.45*0},{0.32*0}) --
        ({4+0.45*3},{0.32*3}) -- ({0+0.45*3},{0.32*3}) -- cycle;
    \draw[dashed, black!25] ({2+0.45*0},{0.32*0}) -- ({2+0.45*3},{0.32*3});
    \draw[dashed, black!25] ({0+0.45*1.5},{0.32*1.5}) -- ({2+0.45*1.5},{0.32*1.5});
    \draw[dashed, black!25] ({2+0.45*2},{0.32*2}) -- ({4+0.45*2},{0.32*2});

    % Prisms, drawn back to front.  Each prism is the graph of the constant
    % prediction over one leaf region: top face at height h, front face
    % (x2 = y0) and left face (x1 = x0) visible.
    \foreach \myx/\myy/\myX/\myY/\myh/\myc in {%
        2/2/4/3/0.1/regiongray,%
        0/1.5/2/3/0.85/mondrianred,%
        2/0/4/2/0.55/mondrianyellow,%
        0/0/2/1.5/0.25/mondrianblue}
    {
        % top face
        \fill[\myc!85]
            ({\myx+0.45*\myy},{0.32*\myy+0.85*\myh}) --
            ({\myX+0.45*\myy},{0.32*\myy+0.85*\myh}) --
            ({\myX+0.45*\myY},{0.32*\myY+0.85*\myh}) --
            ({\myx+0.45*\myY},{0.32*\myY+0.85*\myh}) -- cycle;
        \draw[black!50]
            ({\myx+0.45*\myy},{0.32*\myy+0.85*\myh}) --
            ({\myX+0.45*\myy},{0.32*\myy+0.85*\myh}) --
            ({\myX+0.45*\myY},{0.32*\myY+0.85*\myh}) --
            ({\myx+0.45*\myY},{0.32*\myY+0.85*\myh}) -- cycle;
        % front face (x2 = y0)
        \fill[\myc!60]
            ({\myx+0.45*\myy},{0.32*\myy}) --
            ({\myX+0.45*\myy},{0.32*\myy}) --
            ({\myX+0.45*\myy},{0.32*\myy+0.85*\myh}) --
            ({\myx+0.45*\myy},{0.32*\myy+0.85*\myh}) -- cycle;
        \draw[black!50]
            ({\myx+0.45*\myy},{0.32*\myy}) --
            ({\myX+0.45*\myy},{0.32*\myy}) --
            ({\myX+0.45*\myy},{0.32*\myy+0.85*\myh}) --
            ({\myx+0.45*\myy},{0.32*\myy+0.85*\myh}) -- cycle;
        % left face (x1 = x0)
        \fill[\myc!45]
            ({\myx+0.45*\myy},{0.32*\myy}) --
            ({\myx+0.45*\myY},{0.32*\myY}) --
            ({\myx+0.45*\myY},{0.32*\myY+0.85*\myh}) --
            ({\myx+0.45*\myy},{0.32*\myy+0.85*\myh}) -- cycle;
        \draw[black!50]
            ({\myx+0.45*\myy},{0.32*\myy}) --
            ({\myx+0.45*\myY},{0.32*\myY}) --
            ({\myx+0.45*\myY},{0.32*\myY+0.85*\myh}) --
            ({\myx+0.45*\myy},{0.32*\myy+0.85*\myh}) -- cycle;
        % prediction label at the centre of the top face
        \node[fill=white, fill opacity=0.75, text opacity=1,
              inner sep=1pt, rounded corners=1pt, font=\scriptsize]
            at ({(\myx+\myX)/2+0.45*(\myy+\myY)/2},%
                {0.32*(\myy+\myY)/2+0.85*\myh}) {$c_u=\myh$};
    }

    % Prediction axis, drawn outside the surface.
    \draw[<->] (-0.45,-0.1) -- (-0.45,0.95)
        node[midway, left, font=\small, align=right] {prediction\\$h_T(x)$};
    \node[font=\small] at (2.3,-0.55) {feature space $(x_1,x_2)$};
\end{tikzpicture}
\caption[Piecewise-constant prediction surface of a decision tree.]{The predictor induced by the tree of Figure 9.1 is piecewise constant. Each leaf region $R_u$ is drawn as a prism whose height is the leaf prediction $c_u$; the graph of $h_T$ is the union of the top faces. The surface is flat inside every region and jumps at the partition boundaries, so a tree represents a function with no variation within a leaf and discontinuities across leaves. The figure is the decision-tree analogue of the piecewise-constant surface in Hastie et al. (2009, Figure 9.3), drawn here in the notation of this chapter.}
\label{fig:decision-tree-surface}
\end{figure}
